\section{Scalar lifts on signed arcs of length one and two}

% BEGIN S83L1
\paragraph{S83L1. Two ways of lifting a single coefficient.}
Keep the actual two lattices, surjection $F$, cuts, and constants
$R,\sigma,U$ of U83M1--U83M3. Write $E_2$ for a single integral
constant basis vector at the predecessor, identified by crystalline
Frobenius with its vector at the cut. The coordinate must belong to
$\mathrm{Fil}^{b_u}$ at every up position $u$, with size $b_u$.
Write $P_u$ for the primitive factor there, and put
$D_u=\prod_{u\ {\rm up}}P_u^{b_u}$. The native scalar condition
is $f\in D_u^{-1}A$, $A=O[[v]]$, and for every down position $d$
there are $k_d\ge0$, $x_d\in O$ with
\[
      p^{k_d} f-x_d\in (P_d^{b_d+S}/D_u)A .
\]
This says its first $b_d+S$ Taylor coefficients there agree with a
constant after increasing the coefficient denominator. The cut
idempotent is understood in every scalar.

First, if $D_u f$ is divisible integrally by $P_d^{b_d+S}$ at
each down $d$, then $fE_2\in\mathcal N_2$ of U83M3. Indeed its poles
at the up levels have order at most $b_u$ and the direction is
in the specified filtration there. At the down levels its order is
at least $b_d+S$ and all jumps are at least $-b_d-S$. Hence
$V_2(fE_2)\in L_{2,c}$ by U83M6. U83M1 lifts $p^R$ times this
physical vector in $L_{1,c}$. Every such lattice lift $w_1$ is a
permitted cut input since $\gamma$ acts trivially modulo $v$ on the
invariant Wach lattice (in the full parameter the action is trivial
at zero, and a tame-invariant entry vanishing there is in $vA$).
The image equals $p^R V_2(fE_2)$ and realizes $p^{R+\sigma}f E_2$.
It satisfies the translated-span condition by construction of $V_2$.

Second, suppose an integral predecessor-coordinate vector $C_1$
maps to $p^t E_2$ and belongs, after transport, to
$\mathrm{Fil}^{b_u}$ at all up positions. Form $g=f C_1$.
Then $(\gamma-1)g/v\in\mathcal N_1$. It is integral with
denominator $D_u$: $\gamma$ preserves the denominator ideal, and
the numerator vanishes at zero because $D_u(0)\ne0$ and the value
at zero is unchanged. At an up level its pole order is still at
most $b_u$. At each down level $\gamma$ fixes the scalar constant
$x_d/p^{k_d}$, preserves the Taylor ideal, and $v$ is a unit there.
The difference therefore vanishes to order at least $b_d+S$.
These are exactly the filtration tests of U83M3.
U83M6 gives the source cut vector $V_1(g)$ with
$(\gamma-1)V_1(g)\in vL_{1,c}$.
Its image realizes $p^{t+\sigma}f E_2$. To check this last
assertion, write $M(v)$ for $F$ in predecessor bases and
$H_1,H_2$ for their horizontal matrices normalized at zero.
Uniqueness of normalized horizontal solutions gives
$M(v)H_1(v)=H_2(v)M(0)$ in rational formal series.
Truncate modulo $v^M$, with $M$ the native integer precision,
and apply the arc's Frobenius. This gives equality modulo $v\Lambda^2$
in the translated numerator comparison of U83M6, since
$\varphi^e(v)^M$ is divisible by that modulus. The two physical
frames have integral representatives after $p^\sigma$; the finite
additional clearings in this congruence are allowed by
$\mathrm{RepVV}$. The constant map is $M(0)$, hence sends $C_1$
as asserted. Also $\Lambda^2 V_1(g)\in L_{1,c}$ by the two
height sandwiches, and $\Lambda F V_1(g)$ is in the translated
target span. In both constructions these are genuine
$\mathrm{LiftsAt}$ witnesses with the target correction $z=0$.

For a constant $f=x\in O$ the second argument works for any
integral $C_1$ mapping to $p^t E_2$, because $\gamma g-g=0$
identically. Such a vector always exists with $t=R$: lift $p^R$
times the predecessor basis vector at that cut by U83M1 and
reduce its integral coordinates modulo $v$. Frobenius transports
this equality to the terminal component. Consequently a constant
plus a term of the first kind lifts for the common exponent
$R+\sigma$, by adding the two source vectors. These arguments
apply for any order of the signed letters.
% END S83L1

% BEGIN S83L2
\paragraph{S83L2. The two length-one signs and an integral residue in length two.}
For an up singleton there are no down tests, so the first
construction of S83L1 applies to every $f$. For a down singleton
there is no denominator, and $f\in O[[v]]$. Since $P_0=v+p$,
take $x=f(-p)\in O$. The given condition says $f-x$ vanishes
to order at least $b_0+S$ at $-p$. Division by the distinguished
polynomial $P_0^{b_0+S}$ proves integral divisibility
(the quotient by the ideal is free over $O$). S83L1 now lifts
the constant and remainder. Thus at either length-one sign the
native scalar exponent is $R+\sigma\le6U$. This does not use
the band indicator or an auxiliary rank-one lattice.

At length two let $u$ be the up level, $d$ the down level,
$b_u$ the up size and $y=b_d+S$. Replace the $P_i$ by their
monic distinguished associates when testing divisibility. Write
$Q$ for the associate at $u$ and $P$ for the one at $d$, so
their degrees are $p^u,p^d$ by A82N1. For $f=N/Q^{b_u}$ the
scalar test says $N\equiv\alpha Q^{b_u}\bmod P^y$ over $L$,
with $N\in A$ and $\alpha=x_d/p^{k_d}$. The quotient $A/P^y$
is free over $O$ in its polynomial basis. If
$p^u b_u<p^d y$, the remainder of $Q^{b_u}$ is itself and has
leading coefficient one. Comparison of that coefficient makes
$\alpha$ integral. Thus $f$ splits into an integral constant and
a term satisfying the first construction of S83L1. Both signs
in this strict-degree range again have exponent $R+\sigma$.
% END S83L2

% BEGIN S83L3
\paragraph{S83L3. A source direction in the filtration from a visible coefficient.}
Suppose length two ends at an up letter, with size $a=b_c$.
The previous down letter has size $y-S=b_{c-1}$, and its height
is bounded by $G=w+S$; put $Q=P_0=v+p$ and let $P$ be the
monic degree-$p$ associate of $P_1$. If
\[
                    a>p(y+\eta)+\eta ,
\]
the target test $P^y E_2/Q^a$ is in $\mathcal N_2$ by S83L1.
Use $\Lambda=Q^h P^G$ up to its harmless unit. U83M6 gives
a physical target in $L_{2,c}$ representing $p^\sigma$ times
this test. Lift its $p^R$ multiple into $L_{1,c}$ by U83M1.
There is an integral source numerator $N_1$ coding $p^\sigma$
times this lift in constant coordinates. Indeed
$\Lambda L_{1,c}\subset A\varphi^2 L_{1,c-2}$ and the inverse
horizontal matrix, normalized at zero, has the same degree-by-degree
denominator bound as the direct one: its transpose is the
normalized horizontal solution in the dual basis. U83M2's monomial
calculation applies entry by entry to the truncated inverse too.
The forward and inverse truncated matrices multiply to the identity
modulo $v^M$ before Frobenius. Consequently $N_1/\Lambda$ agrees
with $p^\sigma$ times the source lift through order $qM-h$ at the
up level and $2G-G$ at the down level ($q=p-1$). These errors
lie in the true Taylor lattices. Functoriality in S83L1 therefore
gives, at the up level, equality of target numerators through
degrees strictly below $qM$.

At the down level no jump exceeds $\eta$, so $N_1$ is divisible by
$P^{G-\eta}$ integrally, by distinguished division on each coordinate.
Set $N'_1=N_1/P^{G-\eta}$. At the up level $P$ is a Taylor unit.
All Taylor coefficients of $N'_1$ in degrees $m<h-\eta$ belong to
$\mathrm{Fil}^{\eta+1}=\mathrm{Fil}^a$: remove the pole $Q^{-h}$
in the true Taylor condition and use the gap between the central
window and the cloud. This initial-coefficient assertion survives any
invertible change between local parameters, and we use $Q$ itself.
In the target the corresponding polynomial, up to the scalar
$p^{R+2\sigma}$, is
\[
                        Q^{h-a}P^{y+\eta} E_2 .
\]
Its degree $m=h-a+p(y+\eta)$ is below $h-\eta$ by hypothesis and
its leading coefficient is $E_2$, since $P$ is monic also as a
polynomial in $Q$. It is visible because $qM>2h$. Extract the
coefficient of $Q^m$ in $N'_1$; translation by $-p$ preserves
integrality of all coefficients. This gives a filtered integral
source direction mapping to $p^{R+2\sigma}E_2$. The second
construction in S83L1 gives $E_1=R+3\sigma\le12U$ for every
allowed scalar.

For a terminal down letter instead, write $a=b_{c-1}$ for the
previous up size. Set $I=\lceil s_{c-1}/q\rceil$ where $s_{c-1}$
is the least full Wach $X$-power of R82W1, so $I$ is the native
$\mathrm{AI}(c-1)$. If $a\ge I+\eta+1$, apply the preceding
coefficient argument on the first step only, from $c-2$ to $c-1$.
The parameter is $Q=v+p$, the height $G$, the testing vector
$E_2/Q^a$, and the two direct and inverse truncations have
accuracy $Q^{2G}$ on the numerators. This follows directly from
U83M2 at the first nonthick letter; U83M6's matrix calculation
uses just this first step. The target test lies in its true
filtered lattice by the up-direction condition, transported back
by one crystalline Frobenius. Lift its physical vector multiplied
by $v^I$ on that letter. This is possible since $v$ is $X^q$
times a unit, $qI\ge s_{c-1}$, and averaging over the tame group
of order $p-1$ preserves that image. The target numerator is now
$p^{2\sigma}Q^{G-a}v^I E_2$. Its $Q$-degree is
$G-a+I<G-\eta$ with leading coefficient $p^{2\sigma} E_2$.
The coefficient is thus a filtered integral source direction,
by the same initial-degree test. Transport forward in the definition
of the full arc filtration. This direction lets S83L1 realize every
scalar on the full two-letter arc with $E_1=3\sigma\le9U$.
% END S83L3

% BEGIN S83L4
\paragraph{S83L4. The residue denominator in the other terminal-up case.}
Retain the terminal up orientation and assume $a\le p(y+\eta)+\eta$.
The case $a<py$ has exponent $R+\sigma$ by S83L2. For $py\le a$,
use the actual primitive factor $P=P_1$ and $Q=P_0=v+p$.
Here $P=p\varepsilon+uQ^p$ with $u,\varepsilon$ integral units:
in the tame ring $\varphi(v)=u vQ^{p-1}$ (both sides have the
same cyclotomic factors, by tame descent of the full Frobenius).
Then $P=p+u vQ^{p-1}$ and $\varepsilon=1-uQ^{p-1}$.
Write $a_0=\lfloor a/p\rfloor+1$, $i_0=\min(a_0,y-1)$ and
\[
           d_0=a_0-i_0+v_p\binom{a_0}{i_0}.
\]
For the residue $\alpha$ of S83L2 we claim $p^{d_0}\alpha\in O$.
Indeed $\alpha Q^a$ is integral in $(A/P^y)[1/p]$ by the scalar
test. Multiply it by $u^{a_0}Q^{pa_0-a}$. The resulting integral
class is
\[
   \alpha(P-p\varepsilon)^{a_0}
       =\sum_{i=0}^{i_0}\alpha\binom{a_0}{i}
              (-p)^{a_0-i}\varepsilon^{a_0-i}P^i\pmod {P^y}.
\]
If any displayed scalar coefficient had negative uniformizer
valuation, rescale by a uniformizer power so the smallest becomes
zero. The integral class would reduce to zero modulo the
uniformizer. But over the residue field $P$ has exact series order
$p$ and $\varepsilon$ is a unit. Terms with valuation zero have
distinct leading orders $pi<py$, so their first one survives in the
quotient by $P^y$. This contradiction proves all coefficients
integral, in particular the last one. The argument works over the
original ramified coefficient DVR with $v_p(p)=1$.
Distinguished division now gives a decomposition of $p^{d_0} f$
into the integral constant $p^{d_0}\alpha$ and a scalar with the
full required down zero. S83L1 gives the exponent $R+\sigma+d_0$.

Legendre's formula gives
$v_p\binom{a_0}{i_0}\le\lfloor\log_p(a_0)\rfloor$:
each power $p^k\le a_0$ contributes either zero or one to the
difference of the three factorial valuations. As $py\le a$ and
$y\ge1$, we get
\[
 d_0\le a/p-y+2+\log_p(2h)
      \le\eta+\eta/p+2+\log_p(2h)\le\eta+3U .
\]
This is an actual band place. Put $s=b_{c-1}$, so $y=s+S$.
The adjacent pair consists of the nonthick $s$ and the thick
$a\in[B,B+\eta]$. From $p(s+S)\le a\le p(s+S+\eta)+\eta$
one gets $ps-B\le\eta$ and $B-ps\le pS+p\eta+\eta$.
Thus the original band test holds already with coefficient two
on the expression $p\eta+pS+B/p+s+p\log_p(2h)$, and hence with
the family's $320$. The exponent is at most $9U+\eta I_{\rm band}$.
Combining with S83L2--S83L3 establishes the native scalar comparison
for every terminal-up arc of length two. In the terminal-down
orientation S83L3 supplies the filtered source case $b_{c-1}\ge\mathrm{AI}(c-1)+\eta+1$.
% END S83L4
